% ============================================================================
%  draft_sections.tex  --  paste-ready fragments for Draft v3
%
%  Each block below is labelled with where it goes.  Blocks marked [REPLACE]
%  substitute text already in the draft; blocks marked [NEW] fill a [TODO].
%  Every number quoted here is taken from a cached result in results/ -- the
%  provenance is given in a comment above each block so you can re-check it.
%
%  Still not written here, and why:
%    - Appendix C (reproducibility protocol) -- blocked on SCRTP access
%    - Figure 1 (schematic) -- a drawing, not prose
%    - Appendix D (B corrections) -- untested; your designated first cut
%    - Section 3.5 per-diaphragm disorder -- not yet run
%
%  TWO INCONSISTENCIES IN THE DRAFT THAT I HAVE NOT SILENTLY FIXED -- decide
%  which way you want them, then make them consistent everywhere:
%
%  (1) Which run is "canonical".  Section 2.5 quotes N=200, dt=1e-4, which is the
%      reproducible-result protocol point of Section 3.5.  The dimensional sweeps
%      behind Sections 3.2-3.3 use N=300, dt=5e-4, and the dimensionless sweeps
%      use N=150-200 with T scaled per run.  All are converged, but a reader
%      reproducing Section 3.2 from Section 2.5 will use the wrong one.  Either
%      state the protocol point in 2.5 and note that sweeps size N and T per run,
%      or give a small table.
%
%  (2) Speed extraction.  The draft's 2.5 bullet said "tanh fitting to the front".
%      The code does not do that: it takes the last site below the midpoint,
%      interpolates to sub-site resolution, and fits a straight line to position
%      against time.  The tanh fit exists in the pipeline but is used for the
%      gradient integral I, not for the speed.  My 2.5 below describes the code.
% ============================================================================


% ============================================================================
% [REPLACE]  Abstract -- three numbers change
%   10^-8      -> 10^-14   (results: notebook cell prints 1.0e-14)
%   ~15% / 3 sites -> the real accuracy (results: validity map)
%   X%         -> 9.7%     (results: uq_reproducible_result.ipynb)
% ============================================================================

Transition waves in bistable lattices propagate by converting stored potential
energy into a self-sustaining front, but existing realisations couple their
elements through springs or magnets. We study a one-dimensional chain of bistable
diaphragms coupled instead by a compressible gas obeying Boyle's law --- the
ingredient that makes the array a metafluid rather than a metamaterial.
Non-dimensionalisation reduces the discrete dynamics to four groups
$(\eta,\theta,\Pi,\Omega)$, and simulations at matched groups collapse to within
$10^{-14}$, at integrator round-off. Writing the gas force exactly as a linear
spring stiffness $\eta\theta$ multiplied by a Boyle correction separates the two
roles the fluid plays: $\eta\theta$ sets the front width, which scales as
$\sqrt{\eta\theta}$ and determines whether a travelling profile exists at all,
while the correction contributes a nonlinearity no spring chain can produce. We
map the validity of the energy-balance speed formula, finding agreement to within
$3\%$ wherever the front spans at least three lattice sites and to within $10\%$
down to two, with predictable failure below, where Peierls--Nabarro pinning stalls
the wave entirely. The stall threshold is not a fixed front width but follows
$w_c \propto -\ln\Pi$, the signature of a barrier exponential in width competing
against the energy released per snap. The fluid fingerprint is a
compression/rarefaction speed gap, linear in $\theta$ and reaching $10.8\%$ at
$\theta = 0.4$, identically zero under spring coupling. The measured speed carries
a parametric uncertainty of $9.7\%$ under manufacturing-scale variation in
$\delta$, $A$ and $\alpha$, against a numerical floor of $0.05\%$.


% ============================================================================
% [NEW]  Section 1 -- the three background paragraphs
% ============================================================================

A bistable element stores energy in its higher-energy state and gives it up on
switching. Chain such elements together and a disturbance that switches one can
switch its neighbour, and that neighbour's neighbour in turn, so a front travels
along the chain releasing stored energy as it goes. Nadkarni et al.\
\cite{nadkarni2014dynamics,nadkarni2016universal} showed that the speed of such a
front follows from a balance between the energy released per element and the
energy dissipated per element, and Raney et al.\ \cite{raney2016stable}
demonstrated the waves experimentally in a chain of buckled beams coupled by
elastomeric linkages, where they propagate over long distances without decaying.
Hwang and Arrieta \cite{hwang2018input} applied the same balance to bistable
lattices used for energy harvesting. The formula that comes out of that argument,
quoted as \eqref{eq:arrieta} below, is the one this paper tests; it was derived
for elements coupled by linear springs and magnets, and what happens when the
coupling is a fluid instead is not obvious in advance.

The energy-balance argument assumes a travelling profile exists --- that the front
can be written as a fixed shape moving at constant speed. On a discrete lattice
that assumption has a limit. Puglisi and Truskinovsky \cite{puglisi2000mechanics}
analysed a chain of bistable springs and showed how the discrete problem departs
from its continuum counterpart, and the kinetic relations that govern how fast a
lattice front actually moves were developed by Slepyan \cite{slepyan2002models}
and by Truskinovsky and Vainchtein \cite{truskinovsky2005kinetics}. The obstacle
is the Peierls--Nabarro barrier: a periodic energy landscape the front must climb
to advance by one lattice site, absent from any continuum description and
exponentially small when the front is wide compared with the lattice spacing. The
controlling quantity is therefore not the coupling strength or the drive
separately but the ratio of front width to lattice spacing, and Sections
\ref{sec:validity} and \ref{sec:pinning} are organised around it.

The metafluid setting is more recent. Djellouli et al.\
\cite{djellouli2024metafluid} suspended bistable elastomeric shells in a liquid
and showed that the resulting fluid has a compressibility, an optical response and
a viscosity that can be tuned by design, the shells buckling and unbuckling as the
pressure is cycled. Fluid-driven multistable membranes have been studied in a
related setting by Peretz et al.\ \cite{peretz2020underactuated}, and the
mathematical framework for arrays of subwavelength resonators --- the reduced,
discrete description of a many-element system that this paper also relies on ---
has been developed by Ammari, Davies and Hiltunen \cite{ammari2023functional}.
What has not been studied is a chain of bistable elements in which the coupling
itself is the fluid.

Every one-dimensional bistable chain studied to date couples its elements through
springs or magnets. Here the coupling is a compressible gas obeying Boyle's law,
and this is the whole of the novelty: the coupling force is lopsided in a way no
spring can be, and that asymmetry survives into the wave as a measurable,
direction-dependent speed.


% ============================================================================
% [NEW]  Section 1 -- contributions and roadmap
% ============================================================================

The contributions are as follows. \textbf{D1} We derive the fluid-coupled discrete
model and reduce it to four dimensionless groups, writing the gas force exactly as
a spring stiffness times a Boyle correction (Section \ref{sec:boyle}).
\textbf{D2} We verify the simulator: runs at matched groups but very different
dimensional parameters agree to $10^{-14}$ (Section \ref{sec:verification}).
\textbf{D3} We map where the energy-balance speed formula holds, and show the
controlling variable is the front width measured in lattice sites (Section
\ref{sec:validity}). \textbf{D4} We show that weak coupling pins the wave outright
while weak well asymmetry never does, and that the two are one result seen along
two axes (Section \ref{sec:pinning}). \textbf{D5} We measure the
compression/rarefaction speed asymmetry, the fingerprint of the fluid coupling and
identically absent from any spring chain (Section \ref{sec:fingerprint}). Section
\ref{sec:uq} quantifies the uncertainty on the measured speed; Appendices A and B
give the two derivations the results lean on.


% ============================================================================
% [REPLACE]  Section 3.1 Verification -- clears both tags
%   provenance: dimensionless_model.ipynb "same groups, same wave" cell
%               diaphragm_metafluid.ipynb dt- and N-convergence cells
% ============================================================================

Two physical systems that share the four groups $(\eta,\theta,\Pi,\Omega)$ should
produce the same wave, and this is the sharpest available check on the reduction
from \eqref{eq:eom} to \eqref{eq:nondim}: it exercises the algebra, the
integrator and the boundary treatment at once. We build two systems with
deliberately unlike dimensional parameters --- $(m,k,a,A) = (1, 1, 1, 1)$ and
$(3\times10^{-3},\, 250,\, 0.02,\, 5\times10^{-4})$, a factor of $250$ in
stiffness and $50$ in length scale --- and choose $p_0$, $V_0$, $\delta$ and
$\alpha$ in each so that all four groups match. Integrating both, and separately
integrating the dimensionless equation \eqref{eq:nondim} itself, the three
trajectories agree to $1.0\times10^{-14}$ in the maximum norm over the displacement
field across the whole run. That residual sits at integrator round-off for double
precision, not at a modelling discrepancy, so the reduction is exact to the
accuracy the arithmetic permits.

The timestep is fixed by a Cauchy self-convergence test. The lattice has no closed
form solution to compare against, so instead of measuring $|u_{\text{num}} -
u_{\text{exact}}|$ we halve $\Delta t$ repeatedly and measure the gap between
successive refinements, $|u(\Delta t) - u(\Delta t/2)|$; a sequence converges if
its terms close on each other, whether or not the limit is known. The gap falls
with slope $2$ on log--log, matching the second-order accuracy of velocity Verlet
and confirming the scheme is converging at its design rate rather than being
limited by the damping treatment. Lattice size is checked separately, comparing
matched pairs at $T = 2N$ so that the front traverses the same fraction of the
domain in each and sponge bias does not vary with $N$; the measured speed is flat
in $N$ beyond $N \approx 100$. The residual numerical uncertainty on the wave
speed is $\pm0.0001$ on $0.2206$, or $0.05\%$, and Section \ref{sec:uq} uses this
as the floor beneath the parametric spread.


% ============================================================================
% [REPLACE]  Section 3.2 -- accuracy sentence + the confirmed low-damping claim
%   provenance: results/alpha_low_windows.csv (6 alpha values, 3tau vs 9tau)
%               dimensionless validity map (Omega sweep)
% ============================================================================

Agreement is within $3\%$ wherever the front spans at least three lattice sites,
and within $10\%$ down to a front of about two sites; below that it degrades
quickly and predictably, reaching $23\%$ at $1.3$ sites and $66\%$ at one site,
because the travelling profile $U(\xi)$ assumed by \eqref{eq:arrieta} no longer
exists in any meaningful sense. Since the front width scales as $\sqrt{\kappa}$,
this is a statement about the coupling strength alone.

At low damping the formula overpredicts, and one candidate explanation is a
measurement artefact: the post-snap ringing takes time to die, so a run that is
too short measures a wave that has not yet settled. Appendix \ref{app:wake} shows
that every mode of the wake decays with the same envelope $e^{-\Omega\tilde t/2}$,
so the relevant timescale is $\tilde\tau = 2/\Omega$ and the test is
straightforward: measure the ratio in successive time windows of a single long
run, sizing the lattice to the underdamped speed so the front stays in the bulk.
Across six damping values the ratio is flat to three significant figures between
$3\tilde\tau$ and $9\tilde\tau$ --- $1.803$, $1.525$, $1.269$, $1.125$, $1.045$
and $1.023$ as $\Omega$ increases --- so it does not drift back towards unity once
the ringing has gone, and the overprediction is a genuine steady-state effect
rather than a finite-run-time one. The dimensionless sweep reproduces the same
curve from an independently written integrator, the ratio climbing from $1.017$ to
$2.788$ as $\Omega$ falls from $1.4$ to $0.03$. What the formula neglects in this
regime is the front's own inertia and the lattice waves it radiates, neither of
which fades with time.


% ============================================================================
% [NEW]  Section 3.3 -- replaces the closing [TODO]; the prediction is confirmed
%   provenance: results/pinning_vs_coupling.csv, pinning_finiteT.csv,
%               pinning_boundary.csv;  regenerate with scripts/pinning_boundary.py
% ============================================================================

That reasoning makes a prediction, and it holds. Sweeping $\kappa$ downwards at
the canonical drive, the speed falls steadily and then stops dead between
$\kappa = 0.34$ and $\kappa = 0.30$, a front width between $0.83$ and $0.78$
lattice spacings. The last propagating point still moves at a finite
$v = 4.5\times10^{-3}$, so the arrest is abrupt at this resolution rather than a
smooth approach to zero. Distinguishing a pinned front from a very slow one
matters here, since a front creeping at $v \sim 10^{-4}$ is indistinguishable from
a stationary one in a short run: re-running the first pinned point at $T = 1500$,
$6000$ and $30\,000$ leaves the front at an identical position, $5.587$, with the
fitted speed falling to $-4.5\times10^{-20}$. Over a twenty-fold span of run time
the front does not advance at all.

Repeating the exercise at each of five drives locates the boundary in the
$(\kappa,\Pi)$ plane, and it is not a fixed width. The critical width runs from
$1.10$ lattice spacings at $\Pi = 0.03$ down to $0.38$ at $\Pi = 0.30$, a factor
of $2.9$ across a tenfold change in drive, so ``the front is narrower than a
lattice spacing'' is not the criterion --- at weak drive a $1.1$-spacing front
pins, while at strong drive a $0.38$-spacing front still runs. What does hold is
the form the barrier argument predicts. If the barrier falls exponentially with
width and depinning occurs when the drive matches it, then $w_c$ should be linear
in $\ln\Pi$, and it is:
%
\begin{equation}
  w_c = 0.085 - 0.303\,\ln\Pi, \qquad R^2 = 0.964,
  \label{eq:wc}
\end{equation}
%
against $R^2 = 0.876$ for the best pure power law. Read the other way,
\eqref{eq:wc} says that at fixed low $\kappa$ there is a threshold in $\Pi$ after
all: at $\kappa = 0.32$ the wave pins below $\Pi = 0.115$, and at $\kappa = 0.074$
below $\Pi = 0.30$. The absence of a threshold in the $\Pi$ sweep of the previous
paragraph is therefore a statement about canonical $\kappa$, where the barrier is
negligible, and not a general one. The two sweeps are one result.

It is worth being explicit about why the threshold cannot be a continuum effect.
Equation \eqref{eq:arrieta} gives $v^* = \Delta\tilde\psi/(\Omega I)$; as the
coupling weakens the front steepens, $I$ grows, and $v^*$ approaches zero
asymptotically without ever reaching it at finite drive. A hard threshold ---
$v$ exactly zero at finite $\kappa$ and finite $\Pi$ --- is not available to any
continuum description, so its existence is by itself evidence that the lattice is
doing the work.


% ============================================================================
% [NEW]  Section 2.5 Numerical method
%   NOTE: the draft's bullet said "tanh fitting to the front".  The code does
%   NOT do this -- it uses midpoint crossing with sub-site interpolation, then a
%   linear fit.  Written below to match the code.  The tanh fit exists elsewhere
%   in the pipeline, for the gradient integral I, not for the speed.
% ============================================================================

Equation \eqref{eq:eom} is integrated with velocity Verlet. The damping term makes
the acceleration depend on the velocity, so treating it explicitly would drop the
scheme to first order; instead the velocity update is taken semi-implicitly, with
the damping evaluated at the end of the step,
%
\begin{equation}
  \dot u^{\,n+1} =
  \frac{\dot u^{\,n} + \tfrac12\big(a^n + F^{n+1}/m\big)\Delta t}
       {1 + \tfrac12\,\alpha\,\Delta t/m},
  \label{eq:semiimplicit}
\end{equation}
%
which restores second-order accuracy. The timestep $\Delta t = 10^{-4}$ is set by
the stability of the snapping events rather than by the accuracy of the front: the
front itself is resolved at far coarser steps, but a diaphragm crossing the
spinodal moves quickly enough to destabilise the integration if the step is too
long. Convergence in $\Delta t$ and in $N$ is reported in Section
\ref{sec:verification}.

The domain is $N = 200$ sites with fixed ends. The first five sites are held at
$u_{\text{lo}}$ for the whole run, which acts as a sustained trigger at the
low-energy end and removes any dependence of the measured speed on how the wave
was started; the far end is held at $u_{\text{hi}}$. Both boundaries carry a
sponge layer ten sites deep in which the damping is raised smoothly above its bulk
value on a quadratic taper, adding $4.0$ at the outermost site --- one to two
orders of magnitude above the bulk damping, which is $0.05$ to $0.10$ across the
runs reported here. The taper matters: an abrupt change in damping reflects the
radiation it is supposed to absorb. Its purpose is to absorb the lattice waves shed by the front
before they can return and interfere with it, and all measurements are taken while
the front is in the bulk, well clear of the sponge --- a front that reaches the
sponge is slowed by it and gives a spuriously low speed, which is a failure mode
we check for explicitly rather than assume away.

The front position is located at each saved frame as the last site whose
displacement lies below the midpoint $\tfrac12(u_{\text{lo}} + u_{\text{hi}})$,
refined to sub-site resolution by linear interpolation between that site and its
neighbour. The speed is then the slope of a straight-line fit to front position
against time, taken over the steady portion of the run only; the initial transient,
during which the front is still forming and accelerating onto the fixed point of
Appendix \ref{app:fixedpoint}, is discarded before fitting. Position against time
is genuinely straight once that transient has passed, which is the justification
for the linear fit: over a run long enough for the front to cross $980$ sites, the
fitted speed is constant to four significant figures across successive time
windows, and the energy released by the snapping diaphragms matches the energy
dissipated in each of them, as Appendix \ref{app:fixedpoint} requires.

All simulations are written in Python with NumPy \cite{harris2020numpy}; the
surrogate and sensitivity analysis of Section \ref{sec:uq} use scikit-learn
\cite{pedregosa2011scikit}. Code, cached sweep outputs and the scripts that
regenerate them are in the repository given in Appendix \ref{app:protocol}, which
also gives the hardware and the commands needed to reproduce the canonical result.


% ============================================================================
% [NEW]  Appendix B -- closing paragraph; turns asserted theory into verified
%   provenance: results/mode_decay.csv (3 Omega x 3 modes)
% ============================================================================

Both predictions are checked directly against the lattice. Initialising the chain
in a single normal mode $\epsilon_n = A_0\sin(qn)$ with $q = m\pi/(N-1)$ --- an
exact eigenvector of the discrete Laplacian under fixed ends, so the mode does not
mix --- and fitting the decay of its envelope gives a rate equal to $\Omega/2$ to
three decimal places for every combination of damping and wavenumber tested, with
$R^2 = 1.0000$ on the log-envelope fit. The measured frequencies match
$\sqrt{\omega(q)^2 - \Omega^2/4}$ to the same precision, including the small
downward shift the damping imposes, which accounts for every residual. A single
mode is a pure damped oscillator, exactly as \eqref{eq:dispersion} claims.

The wake behind a real front is not one mode but a superposition of many, and the
consequence is worth stating because it looks at first like a failure. Since every
mode carries the same factor $e^{-\Omega\tilde t/2}$,
%
\begin{equation}
  \epsilon_n(\tilde t) = e^{-\Omega\tilde t/2}
  \sum_q A_q \sin(qn)\,\cos\!\big(\omega_d(q)\,\tilde t + \phi_q\big),
  \label{eq:superposition}
\end{equation}
%
so the signal is exactly an exponential envelope multiplying a quasi-periodic
function. The envelope obeys the theory, but the modes beat against one another,
individual peak heights are irregular, and fitting a decay to a single site over a
short window recovers $\Omega/2$ only unreliably. Measuring one mode at a time
removes the beating; measuring the total wake energy, which sums over sites and
averages the cross terms away, is the other clean route and decays at $\Omega$,
twice the amplitude rate.
